\chapter{Pengaturan Tekanan Darah dan Olahraga}\label{ch:b2:blood-pressure}

Berdirilah dengan cepat setelah berbaring, maka, selama sedetik,
setengah liter \omterm{def:b2:blood-circulation:blood}{darah} mengalir turun ke dalam \omterm{def:b2:blood-circulation:vessels}{vena} tungkai Anda;
\omterm{def:b2:heart:anatomy}{jantungnya}, yang menerima lebih sedikit, memompa lebih sedikit,
tekanan di kepalanya turun, dan otaknya --- yang berada \qty{40}{cm} di
atas \omterm{def:b2:heart:anatomy}{jantungnya} dan tidak dapat menyimpan oksigen --- tinggal sedetik
atau dua detik dari pingsan. Hal itu hampir tidak pernah terjadi,
karena dalam satu denyut sebuah refleks telah mengencangkan \omterm{def:b2:blood-circulation:vessels}{arteriolnya}
dan mempercepat \omterm{def:b2:heart:anatomy}{jantungnya}. Berlarilah mengejar bus, maka ototnya
memerlukan sepuluh kali \omterm{def:b2:blood-circulation:blood}{darah} yang ia punya ketika beristirahat;
ototnya memperolehnya dalam semenit, dan tekanan yang
menggerakkannya ditahan tetap sementara hambatan seluruh tubuhnya turun
dua pertiga. Bab ini membahas pengaturan itu: apa yang menetapkan
tekanannya, bagaimana ia diindra, refleks cepat dan hormon lambat yang
membetulkannya, serta bagaimana sistemnya ditata ulang, alih-alih
dikesampingkan, oleh olahraga.

\section{Apa itu tekanannya}

\begin{theorem}[Persamaan peredarannya]\label{thm:b2:blood-pressure:equation}
\index{tekanan arteri rerata}\index{hambatan tepi total}
\omterm{thm:b2:blood-pressure:equation}{Tekanan arteri rerata} $\bar P$ adalah hasil kali apa yang dimasukkan
\omterm{def:b2:heart:anatomy}{jantungnya} dengan apa yang dikeluarkan pembuluhnya:
\[
\bar P - P_{\text{vena}} = Q\,R, \qquad Q = f\times V_s,
\]
dengan $Q$ adalah \omterm{thm:b2:heart:starling}{curah jantungnya} (laju denyut $f$ kali \omterm{prop:b2:heart:cycle}{volume
sekuncup} $V_s$), $R$ adalah \omterm{thm:b2:blood-pressure:equation}{hambatan tepi totalnya} --- \omterm{def:b2:blood-circulation:vessels}{arteriol} seluruh
organnya yang terpasang sejajar --- dan $P_{\text{vena}}$ adalah
tekanan \omterm{def:b2:blood-circulation:vessels}{venanya}, yang mendekati nol. Ketika beristirahat,
$\qty{5}{L/min}\times\qty{1.08}{mmHg.s/mL}$ memberikan \qty{90}{mmHg}.
Tiap pengaturan tekanannya bekerja pada salah satu dari tiga kuantitas:
lajunya, \omterm{prop:b2:heart:cycle}{volume sekuncupnya} (lewat pengisian dan daya kerut pada
\cref{ch:b2:heart}), atau hambatannya (lewat jari-jari \omterm{def:b2:blood-circulation:vessels}{arteriolnya},
berpangkat empat). Karena organnya terpasang sejajar, aliran ke
masing-masingnya adalah bagian tekanannya dibagi hambatannya sendiri:
melebarkan \omterm{def:b2:blood-circulation:vessels}{arteriol} satu organ menaikkan alirannya tanpa mengubah
aliran organ lainnya, asalkan tekanannya ditahan --- dan menahan
tekanannya itulah tujuan pengaturannya.
\end{theorem}
\begin{proof}
Hukum Ohm bagi sebuah fluida: aliran lewat sebuah hambatan adalah
penurunan tekanan melintasinya dibagi hambatannya (Poiseuille,
\cref{ch:b2:blood-circulation}); bagi hambatan yang terpasang sejajar
alirannya berjumlah dan $1/R = \sum 1/R_i$. Tekanan rerata sepanjang
sebuah denyut kira-kira sama dengan tekanan diastoliknya ditambah
sepertiga \omterm{thm:b2:blood-circulation:windkessel}{tekanan denyutnya}: $80 + 40/3 = \qty{93}{mmHg}$.
\end{proof}

\begin{omfigure}
\begin{tikzpicture}[every node/.style={font=\scriptsize}, >=stealth]
  \node[draw, rounded corners, fill=omThm!15, minimum width=2.4cm, minimum height=1cm, align=center] (heart) at (0,0) {jantung\\$Q = f\,V_s$};
  \draw[very thick, omThm] (1.2,0) -- (2.4,0);
  \draw[very thick, omThm] (2.4,-2.4) -- (2.4,2.4);
  \foreach \y/\org/\r in {2.0/{otak}/{$R_1$}, 1.0/{otot jantung}/{$R_2$}, 0/{usus dan hati}/{$R_3$}, -1.0/{ginjal}/{$R_4$}, -2.0/{otot, kulit}/{$R_5$}} {
    \draw[very thick, omThm] (2.4,\y) -- (3.6,\y);
    \draw[thick, decorate, decoration={zigzag, segment length=4pt, amplitude=2pt}] (3.6,\y) -- (5.0,\y);
    \node[above] at (4.3,\y+0.05) {\org};
    \node[below] at (4.3,\y-0.05) {\r};
    \draw[very thick, omDef] (5.0,\y) -- (6.2,\y);
  }
  \draw[very thick, omDef] (6.2,-2.4) -- (6.2,2.4);
  \draw[very thick, omDef] (6.2,0) -- (7.4,0) |- (0,-3.2) -- (0,-0.5);
  \node[omThm, above] at (1.8,0.05) {$\bar P$};
  \node[omDef, below] at (3.7,-3.2) {vena, $P_{\text{vena}} \approx 0$};
  \node[anchor=west, align=left] at (7.8,1.5) {organ terpasang sejajar:\\$1/R = \sum 1/R_i$\\aliran ke organ $i$: $\bar P/R_i$};
\end{tikzpicture}

{\small Peredarannya sebagai sebuah rangkaian: satu pompa, satu
tekanan, dan hambatan \omterm{def:b2:blood-circulation:vessels}{arteriol} organnya yang terpasang sejajar. Tiap
organ mengambil $\bar P/R_i$; pengaturannya menahan $\bar P$ sementara
tiap $R_i$ disetel menurut permintaannya.}
\end{omfigure}

\begin{proposition}[Mengapa tekanannya harus diatur]\label{prop:b2:blood-pressure:why}
Terlalu rendah, maka otaknya --- yang berada \qty{40}{cm} di atas
\omterm{def:b2:heart:anatomy}{jantungnya} pada orang dewasa yang berdiri, yang berongkos
\qty{31}{mmHg} tinggi tekan hidrostatik --- dan otot \omterm{def:b2:heart:anatomy}{jantungnya}
sendiri kurang teraliri: otaknya gagal dalam hitungan detik. Terlalu
tinggi, maka pembuluhnya rusak: dinding \omterm{def:b2:blood-circulation:vessels}{arterinya} menebal dan
mengeras, glomerulus ginjalnya berparut, retinanya berdarah, dan
\omterm{def:b2:heart:anatomy}{jantungnya} bersusah payah melawan tekanannya sampai ia gagal.
Tekanannya juga harus tetap mantap sementara permintaan organnya
berubah dengan satu orde besaran --- ususnya setelah sebuah santapan,
ototnya ketika berolahraga, kulitnya ketika kepanasan --- dan
sementara volume \omterm{def:b2:blood-circulation:blood}{darahnya} sendiri hilang lewat pendarahan atau
keringat. Dua sistem mengerjakan hal ini: sebuah refleks saraf yang
bekerja dalam satu denyut pada lajunya, \omterm{prop:b2:heart:cycle}{volume sekuncupnya}, dan
hambatannya, serta sistem hormon yang bekerja selama berjam-jam dan
berhari-hari pada hambatannya dan, di atas segalanya, pada volume
\omterm{def:b2:blood-circulation:blood}{darahnya} lewat ginjalnya.
\end{proposition}

\begin{omfigure}
\includegraphics[width=0.6\linewidth]{images/book4/ai/b4-18-blood-pressure-cuff.jpg}

{\small Mengukur tekanannya. Manset dikembungkan di atas tekanan
sistoliknya, sehingga tidak ada \omterm{def:b2:blood-circulation:blood}{darah} yang lewat; ketika ia
dilepaskan, bunyi pertama aliran yang bergolak di bawah stetoskopnya
menandai tekanan sistoliknya, dan lenyapnya bunyi itu, ketika \omterm{def:b2:blood-circulation:vessels}{arterinya}
tetap terbuka sepanjang seluruh daurnya, menandai tekanan diastoliknya.}
\end{omfigure}

\section{Gelung cepatnya: barorefleks}

\begin{definition}[Barorefleks]\label{def:b2:blood-pressure:baroreflex}
\index{baroreseptor}\index{barorefleks}\index{pusat vasomotor}\index{sistem saraf otonom}
Reseptor regangan di dalam dinding sinus karotis (pada percabangan tiap
\omterm{def:b2:blood-circulation:vessels}{arteri} karotis, di bawah rahangnya) dan di dalam lengkung aortanya ---
yaitu \emph{baroreseptor} --- memicu pada laju yang naik bersama
tekanan yang meregangkannya, lebih cepat bagi tekanan yang sedang naik.
Sarafnya mencapai \emph{pusat kardiovaskular} medulanya, yang
mengendalikan kedua cabang keluaran otonomnya: \omterm{prop:b2:heart:control}{saraf vagus}, yang
asetilkolinnya memperlambat \omterm{prop:b2:heart:conduction}{nodus sinoatriumnya}, dan saraf simpatisnya,
yang noradrenalinnya mempercepat nodusnya, menguatkan \omterm{def:b2:heart:anatomy}{biliknya},
menyempitkan \omterm{def:b2:blood-circulation:vessels}{arteriolnya} (menaikkan $R$), dan menyempitkan \omterm{def:b2:blood-circulation:vessels}{venanya}
(memeras \omterm{def:b2:blood-circulation:blood}{darah} yang tersimpan menuju \omterm{def:b2:heart:anatomy}{jantungnya} dan menaikkan
pengisiannya). Kenaikan tekanannya menambah pemicuan baroreseptornya,
yang merangsang \omterm{prop:b2:heart:control}{saraf vagusnya} dan menghambat keluaran simpatisnya:
lajunya turun, \omterm{def:b2:blood-circulation:vessels}{arteriolnya} mengendur, tekanannya turun. Penurunan
tekanan mengerjakan kebalikannya. Gelungnya adalah sebuah \emph{umpan
balik negatif} dengan titik acuan mendekati \qty{95}{mmHg}, tundaan
kira-kira sedetik, dan bati sedemikian sehingga sebuah gangguan
dibetulkan sampai seperlima ukurannya: itulah sebabnya Anda tidak
pingsan ketika berdiri.
\end{definition}
\begin{proof}[Bukti eksperimental]
Hering (1923) menemukan bahwa menekan sinus karotisnya memperlambat
\omterm{def:b2:heart:anatomy}{jantungnya} dan menurunkan tekanannya, dan bahwa memotong saraf
sinusnya meniadakan tanggapannya. Mengaliri sebuah sinus karotis yang
terpencil pada tekanan yang ditetapkan sambil merekam tekanan \omterm{def:b2:blood-circulation:vessels}{arteri}
hewannya memberikan kurva sigmoid refleksnya: tekanan hewannya turun
ketika tekanan sinusnya dinaikkan, dengan curam di sekitar
\qty{100}{mmHg} dan sedikit saja di luar \qtyrange{60}{160}{mmHg}.
Anjing yang saraf sinus maupun saraf aortanya dipotong menjaga tekanan
rerata yang normal selama berhari-hari tetapi tekanan yang berubah
liar dari menit ke menit --- refleksnya adalah sebuah penyangga, bukan
penetap aras jangka panjangnya.
\end{proof}

\begin{omfigure}
\begin{tikzpicture}[every node/.style={font=\scriptsize}, >=stealth]
  \node[draw, rounded corners, fill=omDef!15, align=center] (bp) at (0,0) {tekanan arteri\\$\bar P$};
  \node[draw, rounded corners, fill=omDef!15, align=center] (baro) at (4.2,0) {baroreseptor\\sinus karotis dan lengkung aorta\\laju picu $\uparrow$ bersama $\bar P$};
  \node[draw, rounded corners, fill=omProp!20, align=center] (med) at (8.6,0) {medula:\\pusat kardiovaskular};
  \node[draw, rounded corners, fill=omThm!15, align=center] (vag) at (8.6,-2.2) {vagus $\uparrow$: laju $\downarrow$};
  \node[draw, rounded corners, fill=omThm!15, align=center] (sym) at (4.2,-2.2) {simpatis $\downarrow$: laju, daya kerut,\\tonus arteriol dan vena $\downarrow$};
  \draw[->, thick] (bp) -- (baro); \draw[->, thick] (baro) -- (med);
  \draw[->, thick] (med) -- (vag); \draw[->, thick] (med) -- (sym);
  \draw[->, thick] (sym.west) -| (bp.south) node[pos=0.75, left, align=right] {$Q\downarrow$, $R\downarrow$:\\$\bar P\downarrow$};
  \draw[->, thick] (vag.west) -- (sym.east);
  \node[anchor=north, align=center] at (4.3,-3.2) {kenaikan tekanan dijawab dengan penurunan: umpan balik negatif, tundaan sedetik};
\end{tikzpicture}

{\small \omterm{def:b2:blood-pressure:baroreflex}{Barorefleksnya}. Kenaikan tekanan \omterm{def:b2:blood-circulation:vessels}{arteri} menambah pemicuan
\omterm{def:b2:blood-pressure:baroreflex}{baroreseptornya}, yang mendorong medulanya untuk memperlambat \omterm{def:b2:heart:anatomy}{jantungnya}
dan mengendurkan pembuluhnya; penurunan mengerjakan kebalikannya.}
\end{omfigure}

\begin{theorem}[Bati sebuah gelung umpan balik negatif]\label{thm:b2:blood-pressure:gain}
\index{bati umpan balik}
Jika sebuah gangguan akan, tanpa pengaturan, mengubah tekanannya
sebesar $\Delta P_0$, dan gelungnya menanggapi tiap simpangan $\Delta P$ dengan pembetulan sebesar $-G\,\Delta P$ (yaitu \emph{bati gelung
terbuka} $G$), simpangan yang tersisa setelah gelungnya bekerja adalah
\[
\Delta P = \frac{\Delta P_0}{1 + G}.
\]
\omterm{def:b2:blood-pressure:baroreflex}{Barorefleksnya} memiliki $G \approx 4$: berdiri, yang akan menurunkan
tekanan di \omterm{def:b2:heart:anatomy}{jantungnya} sekitar \qty{40}{mmHg}, menurunkannya
\qty{8}{mmHg}; sebuah pendarahan yang akan menyetengahkan tekanannya
menurunkannya sepersepuluh --- sampai kehilangannya melampaui apa yang
dapat disembunyikan pembuluh yang lebih kencang dan \omterm{def:b2:heart:anatomy}{jantung} yang lebih
cepat. Refleksnya tidak meniadakan sebuah gangguan; ia membaginya
dengan lima. Dan ia \emph{menyetel ulang}: jika ditahan pada tekanan
yang baru selama berjam-jam, \omterm{def:b2:blood-pressure:baroreflex}{baroreseptornya} menyesuaikan diri lalu
membela aras yang baru itu, dan itulah sebabnya refleksnya tidak dapat
menyembuhkan hipertensi dan sebabnya aras jangka panjangnya ditetapkan
di tempat lain.
\end{theorem}
\begin{proof}
Simpangan akhirnya adalah gangguannya ditambah pembetulannya:
$\Delta P = \Delta P_0 - G\,\Delta P$, sehingga $\Delta P(1 + G) =
\Delta P_0$. Dengan $G = 4$, $\Delta P = \Delta P_0/5$.
\end{proof}

\begin{omfigure}
\begin{tikzpicture}
\begin{axis}[
  width=0.7\linewidth, height=5.4cm,
  axis x line=bottom, axis y line=left,
  xmin=40, xmax=180, ymin=40, ymax=160,
  xlabel={tekanan di dalam sinus karotis (\unit{mmHg})}, ylabel={tekanan arteri yang dihasilkan (\unit{mmHg})},
  xlabel style={font=\small, at={(axis description cs:0.5,-0.13)}, anchor=north},
  ylabel style={font=\small},
  tick label style={font=\small}, xtick={60,100,140}, ytick={60,100,140},
  clip=false,
]
\addplot[very thick, omDef, domain=40:180, samples=80] {100 + 50*tanh((100-x)/25)};
\addplot[thick, dashed, gray, domain=40:180, samples=2] {x};
\fill[omDef] (axis cs:100,100) circle (2pt); \draw[gray] (axis cs:101,100) -- (axis cs:124,110); \node[font=\scriptsize, anchor=west] at (axis cs:125,110) {titik acuan};
\node[font=\scriptsize, anchor=west, align=left] at (axis cs:46,112) {kemiringan $-G$\\dekat titik acuannya};
\end{axis}
\end{tikzpicture}

{\small Kurva \omterm{def:b2:blood-pressure:baroreflex}{barorefleksnya}: tekanan yang dicapai tubuhnya ketika
sinus yang terpencil ditahan pada sebuah tekanan tertentu.
Kemiringannya pada titik acuannya adalah batinya; di luar
\qtyrange{60}{160}{mmHg} refleksnya jenuh.}
\end{omfigure}

\section{Gelung lambatnya: hormon dan ginjalnya}

\begin{proposition}[Renin, angiotensin, aldosteron, dan volumenya]\label{prop:b2:blood-pressure:raas}
\index{sistem renin--angiotensin}\index{aldosteron}\index{hormon antidiuretik}\index{peptida natriuretik}
Selama berjam-jam dan berhari-hari tekanannya ditetapkan oleh
\emph{volume \omterm{def:b2:blood-circulation:blood}{darahnya}}, dan volumenya oleh ginjalnya. Ketika tekanannya
atau natrium yang diantarkan ke ginjalnya turun, atau saraf simpatisnya
memicu, sel \omterm{def:b2:blood-circulation:vessels}{arteriol} ginjalnya melepaskan enzim \emph{renin} ke dalam
\omterm{def:b2:blood-circulation:blood}{darahnya}; renin memotong sebuah protein \omterm{def:b2:blood-circulation:blood}{plasma} menjadi
\emph{angiotensin I}, yang oleh sebuah enzim \omterm{def:b2:blood-circulation:vessels}{kapiler} parunya diubah
menjadi \emph{angiotensin II} --- penyempit pembuluh yang paling
dahsyat di dalam tubuhnya, yang menaikkan $R$ dalam hitungan menit, dan
sebuah hormon yang membuat korteks adrenalnya menyekresikan
\emph{aldosteron}, yang membuat ginjalnya menahan natrium, dan
bersamanya air, selama hari-hari berikutnya. \emph{\omterm{prop:b2:blood-pressure:raas}{Hormon antidiuretik}}
hipofisisnya (vasopresin), yang dilepaskan ketika \omterm{def:b2:blood-circulation:blood}{darahnya} menjadi
pekat atau volumenya turun, membuat ginjalnya menahan air secara
langsung dan juga menyempitkan pembuluhnya. Dan ketika \omterm{def:b2:heart:anatomy}{serambinya}
teregang berlebihan oleh volume yang terlalu banyak, \omterm{def:b2:heart:anatomy}{serambinya}
menyekresikan \emph{\omterm{prop:b2:blood-pressure:raas}{peptida natriuretik}}, yang membuat ginjalnya
membuang natrium dan air. Karena itu ginjalnya adalah wasit
terakhirnya: ia dapat membuang atau menahan berliter-liter, dan sebuah
tekanan yang bertahan di atas titik acuan ginjalnya dibetulkan, selama
berhari-hari, oleh kehilangan volumenya --- kecuali jika titik acuan
ginjalnya itu sendiri dinaikkan, dan itulah kebanyakan hipertensi.
(Mesin ginjalnya sendiri adalah pokok bahasan jilid Tahun 3.)
\end{proposition}
\begin{proof}[Bukti eksperimental]
Goldblatt (1934) menyempitkan \omterm{def:b2:blood-circulation:vessels}{arteri} menuju salah satu ginjal seekor
anjing: hewannya menjadi hipertensif dalam hitungan hari lalu tetap
begitu; ginjal yang kurang teraliri itu ternyata melepaskan renin, dan
menyingkirkannya menyembuhkan hewannya. Obat yang menghalangi
pengubahan angiotensinnya (penghambat ACE) atau reseptornya menurunkan
tekanan kebanyakan penderita hipertensi, dan sebuah ginjal yang
dicangkokkan dari galur tikus yang hipertensif membuat seekor penerima
yang normal menjadi hipertensif.
\end{proof}

\begin{omfigure}
\begin{tikzpicture}[every node/.style={font=\scriptsize}, >=stealth]
  \node[draw, rounded corners, fill=omDef!15, align=center] (fall) at (0,0) {tekanan atau volume turun;\\saraf simpatis memicu};
  \node[draw, rounded corners, fill=omProp!20, align=center] (renin) at (4.4,0) {ginjal melepaskan\\renin};
  \node[draw, rounded corners, fill=omProp!20, align=center] (ang) at (8.6,0) {angiotensin II\\(dibuat di dalam darah dan paru)};
  \node[draw, rounded corners, fill=omThm!15, align=center] (vc) at (8.6,-2.0) {penyempitan pembuluh:\\$R\uparrow$ dalam hitungan menit};
  \node[draw, rounded corners, fill=omThm!15, align=center] (ald) at (4.4,-2.0) {aldosteron (adrenal):\\Na$^+$ dan air ditahan\\selama berhari-hari --- volume $\uparrow$};
  \draw[->, thick] (fall) -- (renin); \draw[->, thick] (renin) -- (ang);
  \draw[->, thick] (ang) -- (vc); \draw[->, thick] (ang) -- (ald);
  \draw[->, thick] (ald.west) -| (fall.south) node[pos=0.85, left, align=right] {tekanan\\pulih};
  \draw[->, thick] (vc.west) -- (ald.east);
  \node[draw, rounded corners, fill=omExo!25, align=center] (anp) at (0,-2.0) {serambi teregang berlebihan:\\peptida natriuretik,\\garam dan air dibuang};
  \draw[->, thick, dashed] (anp.north) -- (fall.south);
\end{tikzpicture}

{\small Gelung lambatnya. Penurunan tekanan memicu renin, angiotensin
II, dan aldosteron: penyempitan dalam hitungan menit, penahanan garam
dan air selama berhari-hari. Volume yang terlalu banyak dijawab oleh
\omterm{prop:b2:blood-pressure:raas}{peptida natriuretik} \omterm{def:b2:heart:anatomy}{serambinya}.}
\end{omfigure}

\begin{example}[Sebuah pendarahan]\label{ex:b2:blood-pressure:haemorrhage}
Seorang pendonor memberikan \qty{500}{mL}, yaitu sepersepuluh \omterm{def:b2:blood-circulation:blood}{darahnya}.
Dalam hitungan detik \omterm{def:b2:blood-pressure:baroreflex}{barorefleksnya} menyempitkan \omterm{def:b2:blood-circulation:vessels}{arteriol} dan \omterm{def:b2:blood-circulation:vessels}{venanya}
lalu mempercepat \omterm{def:b2:heart:anatomy}{jantungnya}: tekanannya nyaris tidak bergerak, kulitnya
menjadi pucat dan dingin (\omterm{def:b2:blood-circulation:vessels}{arteriolnya} tertutup), dan denyut nadinya
lebih cepat. Dalam hitungan menit tekanan \omterm{def:b2:blood-circulation:vessels}{kapilernya} yang menurun
membiarkan cairan antarselnya berpindah ke dalam pembuluhnya
(perimbangan Starling pada \cref{ch:b2:blood-circulation}), sehingga
memulihkan separuh volumenya pada jam berikutnya; renin, angiotensin,
dan \omterm{prop:b2:blood-pressure:raas}{hormon antidiuretik} menyusutkan air seninya menjadi setetes-setetes,
lalu rasa haus menyusul; selama berhari-hari ginjalnya menahan garam
dan air lalu volume \omterm{def:b2:blood-circulation:blood}{plasmanya} kembali, dan selama berminggu-minggu
sumsumnya mengganti \omterm{def:b2:blood-circulation:blood}{sel darah merahnya}. Kehilangan \qty{2}{L}
mendahului semuanya itu: refleksnya jenuh, tekanannya runtuh, jaringan
yang kurang teraliri melepaskan laktat, dan \omterm{def:b2:blood-circulation:vessels}{kapilernya} bocor --- syok,
yang hanya transfusi dapat memulangkan pasiennya.
\end{example}

\section{Olahraga: sistemnya ditata ulang}

\begin{proposition}[Olahraga]\label{prop:b2:blood-pressure:exercise}
\index{faal olahraga}\index{aliran darah otot}\index{curah jantung!olahraga}
Pada olahraga yang berat pemakaian oksigen otot rangkanya naik dua
puluh kali lipat, dan aliran \omterm{def:b2:blood-circulation:blood}{darahnya} dari \qty{1}{L/min} menjadi
\qty{20}{L/min}. Tiga hal terjadi sekaligus. Secara setempat,
metabolit otot yang bekerja --- $\mathrm{CO_2}$, laktat, adenosin,
kalium, panas, oksigen yang rendah --- mengendurkan \omterm{def:b2:blood-circulation:vessels}{arteriolnya}
(\emph{pelebaran pembuluh metabolis}) lalu membuka \omterm{def:b2:blood-circulation:vessels}{kapilernya},
sehingga hambatan ototnya terpangkas menjadi sepersepuluh. Secara
memusat, sebuah perintah dari korteks motoriknya kepada medulanya,
yang diperkuat isyarat dari reseptor ototnya sendiri, menyetel ulang
\omterm{def:b2:blood-pressure:baroreflex}{barorefleksnya} kepada titik acuan yang lebih tinggi lalu mendorong
keluaran simpatisnya: lajunya menjadi 180, daya kerutnya naik, \omterm{def:b2:blood-circulation:vessels}{arteriol}
usus, ginjal, dan otot yang beristirahat menyempit, dan \omterm{def:b2:blood-circulation:vessels}{venanya}
diperas; pompa ototnya dan pernapasan yang dalam mengembalikan \omterm{def:b2:blood-circulation:blood}{darahnya}
lebih cepat. Hasilnya adalah keluaran \qty{25}{L/min} dengan hambatan
total sepertiga nilai istirahatnya, tekanan sistolik \qty{150}{mmHg},
dan tekanan diastolik yang tidak lebih tinggi daripada ketika
beristirahat; empat perlima alirannya pergi ke ototnya, bagian otaknya
tak berubah secara mutlak, bagian ususnya menyetengah, dan bagian
kulitnya naik untuk membuang panasnya. Refleksnya tidak dikesampingkan
melainkan disetel ulang: ia membela tekanan yang lebih tinggi sementara
pembuluhnya membuka kembali tubuhnya.
\end{proposition}

\begin{omfigure}
\begin{tikzpicture}
\begin{axis}[
  ybar, bar width=9pt,
  width=0.9\linewidth, height=6cm,
  ymin=0, ymax=22, ytick={0,5,10,15,20},
  ylabel={aliran darah (\unit{L/min})}, ylabel style={font=\small},
  symbolic x coords={otak, jantung, usus dan hati, ginjal, otot, kulit, lainnya},
  xtick=data, x tick label style={font=\scriptsize, rotate=30, anchor=north east},
  tick label style={font=\small},
  legend style={font=\scriptsize, at={(0.03,0.97)}, anchor=north west, draw=none},
  enlarge x limits=0.12,
  clip=false,
]
\addplot[fill=omDef!50, draw=omDef] coordinates {(otak,0.75) (jantung,0.25) (usus dan hati,1.4) (ginjal,1.1) (otot,1.0) (kulit,0.3) (lainnya,0.2)};
\addlegendentry{istirahat, \qty{5}{L/min}}
\addplot[fill=omThm!50, draw=omThm] coordinates {(otak,0.75) (jantung,1.0) (usus dan hati,0.6) (ginjal,0.6) (otot,20) (kulit,1.6) (lainnya,0.45)};
\addlegendentry{olahraga berat, \qty{25}{L/min}}
\end{axis}
\end{tikzpicture}

{\small Ke mana \omterm{def:b2:blood-circulation:blood}{darahnya} pergi. Pada olahraga yang berat keluarannya
naik lima kali lipat lalu dibagikan ulang: dua puluh kali lipat ke
ototnya, lebih banyak ke \omterm{def:b2:heart:anatomy}{jantung} dan kulitnya, lebih sedikit ke usus
dan ginjalnya, dan sama ke otaknya.}
\end{omfigure}

\begin{theorem}[Pengantaran oksigen dan asas Fick]\label{thm:b2:blood-pressure:fick}
\index{asas Fick}\index{VO2max@$\dot V_{\mathrm{O_2}\max}$}
Oksigen yang dihabiskan sebuah tubuh tiap menit sama dengan \omterm{thm:b2:heart:starling}{curah
jantungnya} kali selisih kandungan oksigen antara \omterm{def:b2:blood-circulation:blood}{darah} \omterm{def:b2:blood-circulation:vessels}{arteri} dan \omterm{def:b2:blood-circulation:blood}{darah}
\omterm{def:b2:blood-circulation:vessels}{vena} campurannya:
\[
\dot V_{\mathrm{O_2}} = Q\,(C_a - C_v).
\]
Ketika beristirahat, $\qty{5}{L/min}\times(200 - 150)\,\unit{mL/L} = \qty{250}{mL/min}$;
pada olahraga maksimal, $\qty{25}{L/min}\times(200 - 40) =
\qty{4000}{mL/min}$, yaitu kenaikan enam belas kali lipat dari kenaikan
aliran lima kali lipat dan kenaikan penyaripatian tiga kali lipat. Laju
maksimal ini, $\dot V_{\mathrm{O_2}\max}$, adalah ukuran baku bagi
kebugaran daya tahan, dan langit-langitnya adalah \omterm{def:b2:heart:anatomy}{jantungnya}: ototnya
masih dapat menyaripatikan lebih banyak dan paru-parunya masih dapat
memuati lebih banyak, tetapi keluarannya tidak dapat naik lagi.
Latihan menaikkan \omterm{prop:b2:heart:cycle}{volume sekuncupnya} (\omterm{def:b2:heart:anatomy}{bilik} yang lebih besar dan lebih
lentur, volume \omterm{def:b2:blood-circulation:blood}{darah} yang lebih banyak) dan dengan demikian
keluarannya --- laju istirahat 45 seorang atlet adalah tanda \omterm{prop:b2:heart:cycle}{volume
sekuncup} \qty{110}{mL} --- serta menambahkan \omterm{def:b2:blood-circulation:vessels}{kapiler} dan mitokondria
kepada ototnya sehingga ototnya menyaripatikan lebih banyak.
Sebaliknya, \omterm{thm:b2:blood-pressure:fick}{asas Fick} adalah cara keluarannya sendiri diukur pada
seorang pasien, dari oksigen yang dihabiskan dan kedua contoh \omterm{def:b2:blood-circulation:blood}{darahnya}.
\end{theorem}
\begin{proof}
Tiap liter \omterm{def:b2:blood-circulation:blood}{darah} yang melewati jaringannya meninggalkan $C_a - C_v$
mililiter oksigen di belakangnya, dan $Q$ liter lewat tiap menit; pada
keadaan tunak itulah yang diserap paru-parunya dan dihabiskan tubuhnya.
Menyelesaikannya bagi $Q$ memberikan pengukurannya.
\end{proof}

\begin{example}[Berdiri, pingsan, dan sang antariksawan]\label{ex:b2:blood-pressure:standing}
Ketika berdiri, \qty{500}{mL} \omterm{def:b2:blood-circulation:blood}{darah} berpindah ke dalam \omterm{def:b2:blood-circulation:vessels}{vena} tungkainya,
kembalian \omterm{def:b2:blood-circulation:vessels}{vena} dan \omterm{prop:b2:heart:cycle}{volume sekuncupnya} turun sepertiga, dan tekanan di
sinus karotisnya merosot --- \omterm{def:b2:blood-pressure:baroreflex}{barorefleksnya} menjawab dalam satu denyut
dengan \omterm{def:b2:heart:anatomy}{jantung} yang lebih cepat dan pembuluh yang lebih kencang, dan
tekanan di kepalanya pulih dalam sepuluh detik. Seorang prajurit yang
berdiri tak bergerak pada sebuah lapangan upacara yang panas kehilangan
pompa ototnya, menggenangkan \omterm{def:b2:blood-circulation:blood}{darah} di dalam \omterm{def:b2:blood-circulation:vessels}{vena} kulitnya yang melebar,
lalu pingsan: berbaring datar menyembuhkannya seketika. Seorang
antariksawan setelah berbulan-bulan tanpa gravitasi telah kehilangan
satu liter volume \omterm{def:b2:blood-circulation:blood}{darahnya} (ginjalnya membuang apa yang ia baca sebagai
kelebihan ketika \omterm{def:b2:blood-circulation:blood}{darahnya} menggenang di dadanya), dan ketika mendarat
refleksnya, yang tak terlatih, tidak sanggup menjaganya tetap tegak.
Dan jerapah, yang kepalanya berada dua meter di atas \omterm{def:b2:heart:anatomy}{jantungnya},
menjalankan tekanan rerata \qty{200}{mmHg} dan mengenakan kaus kaki
penekan dari kulit --- hidrostatika pada
\cref{ch:b2:blood-circulation} menetapkan syaratnya, dan refleks pada
bab ini memenuhinya.
\end{example}

\section{Latihan}

\begin{exercise}[$\star$]\label{exo:b2:blood-pressure:1}
Tulislah persamaan yang menghubungkan tekanan rerata, \omterm{thm:b2:heart:starling}{curah jantung},
dan hambatan tepinya, lalu sebutkan ketiga kuantitas yang dapat diubah
tubuhnya beserta organ atau jaringan yang mengubah masing-masingnya.
\end{exercise}

\begin{exercise}[$\star$]\label{exo:b2:blood-pressure:2}
Perikanlah \omterm{def:b2:blood-pressure:baroreflex}{barorefleksnya}: pengindra, pusat, efektor, tanda umpan
baliknya, dan tundaannya. Apa yang terjadi pada tekanannya dari menit
ke menit pada seekor hewan yang saraf \omterm{def:b2:blood-pressure:baroreflex}{baroreseptornya} dipotong?
\end{exercise}

\begin{exercise}[$\star$]\label{exo:b2:blood-pressure:3}
Berikan urutan renin $\to$ angiotensin $\to$ aldosteron, beserta organ
yang membuat masing-masingnya dan apa yang dikerjakan masing-masingnya,
lalu katakan apa yang memulainya.
\end{exercise}

\begin{exercise}[$\star$]\label{exo:b2:blood-pressure:4}
Sebutkan ketiga mekanisme yang melipatduapuluhkan aliran \omterm{def:b2:blood-circulation:blood}{darah} sebuah
otot yang bekerja, dan katakan organ mana yang menyerahkan alirannya
agar hal itu mungkin.
\end{exercise}

\begin{exercise}[$\star\star$]\label{exo:b2:blood-pressure:5}
Hitunglah tekanan reratanya bagi $Q = \qty{5}{L/min}$ dan $R =
\qty{1.2}{mmHg.s/mL}$; lalu bagi $Q = \qty{20}{L/min}$ dan $R =
\qty{0.35}{mmHg.s/mL}$. Pada keadaan yang kedua, dengan faktor berapa
jari-jari \omterm{def:b2:blood-circulation:vessels}{arteriol} reratanya telah berubah?
\end{exercise}

\begin{exercise}[$\star\star$]\label{exo:b2:blood-pressure:6}
Sebuah gangguan akan menurunkan tekanannya sebesar \qty{30}{mmHg}.
Hitunglah penurunan sisanya bagi bati refleks 2, 4, dan 8. Sebuah obat
menyetengahkan bati \omterm{def:b2:blood-pressure:baroreflex}{barorefleks} seorang pasien: apa yang ia rasakan
ketika berdiri?
\end{exercise}

\begin{exercise}[$\star\star$]\label{exo:b2:blood-pressure:7}
Rapat massa \omterm{def:b2:blood-circulation:blood}{darah} \qty{1060}{kg/m^3}; otaknya berada \qty{40}{cm} di
atas \omterm{def:b2:heart:anatomy}{jantungnya} pada orang dewasa yang berdiri dan kakinya
\qty{120}{cm} di bawahnya. Hitunglah beda tekanan hidrostatiknya dalam
mmHg, dan tekanan \omterm{def:b2:blood-circulation:vessels}{arteri} di otaknya dan di pergelangan kakinya ketika
tekanan di \omterm{def:b2:heart:anatomy}{jantungnya} \qty{95}{mmHg}. Apa yang dibutuhkan seekor
jerapah yang kepalanya \qty{2.5}{m} di atas \omterm{def:b2:heart:anatomy}{jantungnya}?
\end{exercise}

\begin{exercise}[$\star\star$]\label{exo:b2:blood-pressure:8}
Dengan memakai \omterm{thm:b2:blood-pressure:fick}{asas Fick}, hitunglah \omterm{thm:b2:heart:starling}{curah jantung} seorang pasien yang
menghabiskan \qty{280}{mL} oksigen tiap menit dengan kandungan \omterm{def:b2:blood-circulation:vessels}{arteri}
\qty{190}{mL/L} dan kandungan \omterm{def:b2:blood-circulation:vessels}{vena} campuran \qty{130}{mL/L}. Seorang
atlet menghabiskan \qty{5}{L/min} pada maksimumnya dengan kandungan 200
dan \qty{30}{mL/L}: hitunglah keluarannya dan, pada laju 190, \omterm{prop:b2:heart:cycle}{volume
sekuncupnya}.
\end{exercise}

\begin{exercise}[$\star\star$]\label{exo:b2:blood-pressure:9}
Sebuah pendarahan \qty{800}{mL} akan, tanpa pengaturan, menurunkan
tekanannya sebesar \qty{35}{mmHg}. Berikan tekanannya setelah
\omterm{def:b2:blood-pressure:baroreflex}{barorefleksnya} ($G = 4$), sebutkan keempat hal yang memulihkan
volumenya beserta skala waktunya, dan jelaskan kulitnya yang pucat dan
dingin.
\end{exercise}

\begin{exercise}[$\star\star\star$]\label{exo:b2:blood-pressure:10}
Otot ketika beristirahat memiliki $R = \qty{5}{mmHg.s/mL}$ dan menerima
\qty{1}{L/min}; ketika berolahraga hambatannya turun menjadi
sepersepuluh. Jika tekanannya ditahan pada \qty{95}{mmHg} dan tidak ada
lainnya yang berubah, aliran berapa yang akan ditarik ototnya, dan
keluaran berapa yang akan dibutuhkan \omterm{def:b2:heart:anatomy}{jantungnya}? Jika sebagai gantinya
keluaran \omterm{def:b2:heart:anatomy}{jantungnya} dipancangkan pada \qty{5}{L/min}, tekanannya akan
turun menjadi berapa? Jelaskan mengapa tubuhnya tidak mengerjakan
keduanya dan apa yang ia kerjakan.
\end{exercise}

\begin{exercise}[$\star\star\star$]\label{exo:b2:blood-pressure:11}
Goldblatt menyempitkan salah satu \omterm{def:b2:blood-circulation:vessels}{arteri} ginjalnya; anjingnya menjadi
hipertensif secara tetap, dengan tanggapan \omterm{def:b2:blood-pressure:baroreflex}{barorefleks} yang normal.
Jelaskan (a) mengapa ginjalnya menaikkan tekanannya, (b) mengapa
\omterm{def:b2:blood-pressure:baroreflex}{barorefleksnya} tidak mencegahnya, (c) mengapa menyingkirkan ginjalnya
menyembuhkannya, (d) apa yang akan dikerjakan sebuah penghambat ACE.
\end{exercise}

\begin{exercise}[$\star\star\star$]\label{exo:b2:blood-pressure:12}
``\omterm{def:b2:blood-pressure:baroreflex}{Barorefleksnya} adalah sebuah penyangga; ginjalnya adalah sebuah
termostat.'' Bahaslah skala waktu, bati, dan titik acuan kedua
sistemnya, dan apa yang dapat dan tidak dapat dikerjakan
masing-masingnya terhadap tekanan yang naik secara menahun.
\end{exercise}

\section{Soal: Berdiri, Berdarah, Berlari}

\begin{problem}[{Soal akhir pekan --- satu peredaran dibawa melewati
tiga tantangan, dengan tekanan, aliran, pembetulan refleks, dan
pengantaran oksigennya dihitung, berakhir pada penurunan sisa ketika
berdiri, tekanan setelah sebuah pendarahan, serta keluaran dan serapan
oksigennya ketika berolahraga}]
\label{pb:b2:blood-pressure:1}
Data ketika beristirahat: $Q = \qty{5}{L/min}$, laju 70, $\bar P = \qty{95}{mmHg}$,
$P_{\text{vena}} = 0$, \omterm{def:b2:blood-circulation:blood}{darah} \qty{5}{L}, bati \omterm{def:b2:blood-pressure:baroreflex}{barorefleks} $G = 4$,
oksigen \omterm{def:b2:blood-circulation:vessels}{arteri} \qty{200}{mL/L}, oksigen \omterm{def:b2:blood-circulation:vessels}{vena} \qty{150}{mL/L}. Hambatan
organ ketika beristirahat (\unit{mmHg.s/mL}): otak 7.6, \omterm{def:b2:heart:anatomy}{jantung} 22.8,
usus dan hati 4.1, ginjal 5.2, otot 5.7, kulit 19, lainnya 28.5. Rapat
massa \omterm{def:b2:blood-circulation:blood}{darah} \qty{1060}{kg/m^3}; $\qty{1}{mmHg} = \qty{133}{Pa}$.

\medskip
\noindent\textbf{Bagian I --- Ketika beristirahat.}
\begin{enumerate}
  \item Hitunglah \omterm{thm:b2:blood-pressure:equation}{hambatan tepi totalnya} dari hambatan organnya yang
        terpasang sejajar, lalu periksalah ia terhadap $\bar P/Q$.
  \item Hitunglah aliran ke tiap organnya dan bagian keluaran yang ia
        terima.
  \item Hitunglah \omterm{prop:b2:heart:cycle}{volume sekuncupnya} dan pemakaian oksigennya.
  \item Otaknya berbobot \qty{1.4}{kg} dan ginjalnya \qty{0.3}{kg}.
        Hitunglah alirannya tiap kilogram lalu berilah ulasan.
  \item Mengapa tekanannya, dan bukan alirannya, yang harus menjadi
        peubah yang diatur?
  \item Otot \omterm{def:b2:heart:anatomy}{jantungnya} menerima \qty{250}{mL/min} dan menyaripatikan
        \qty{70}{\%} oksigennya; otot ketika beristirahat
        menyaripatikan \qty{25}{\%}. Hitunglah oksigen yang dipakai
        masing-masingnya tiap menit dan jelaskan mengapa satu-satunya
        cadangan \omterm{def:b2:heart:anatomy}{jantungnya} adalah aliran yang lebih banyak.
\end{enumerate}

\medskip
\noindent\textbf{Bagian II --- Berdiri.}
\begin{enumerate}[resume]
  \item Ketika berdiri, \qty{500}{mL} berpindah ke dalam \omterm{def:b2:blood-circulation:vessels}{vena}
        tungkainya dan \omterm{prop:b2:heart:cycle}{volume sekuncupnya} turun \qty{30}{\%}.
        Hitunglah penurunan $\bar P$ tanpa pengaturan (hambatannya tak
        berubah).
  \item Hitunglah penurunan sisanya setelah \omterm{def:b2:blood-pressure:baroreflex}{barorefleksnya}.
  \item Refleksnya mencapai hal ini dengan menaikkan lajunya menjadi
        85 dan menaikkan hambatannya. Hitunglah hambatan baru yang
        dibutuhkan.
  \item Hitunglah beda tekanan hidrostatik antara \omterm{def:b2:heart:anatomy}{jantungnya} dan
        otaknya yang \qty{40}{cm} lebih tinggi, serta tekanan \omterm{def:b2:blood-circulation:vessels}{arteri}
        otaknya sebelum dan sesudah refleksnya.
  \item Seorang pasien yang memakai obat penghalang saraf simpatis
        memiliki $G = 1$. Hitunglah penurunan sisanya dan tekanan
        otaknya, lalu katakan apa yang ia alami.
  \item Jelaskan mengapa membaringkan pasien yang pingsan secara datar
        memulihkan kesadarannya dalam hitungan detik.
\end{enumerate}

\medskip
\noindent\textbf{Bagian III --- Berdarah.}
\begin{enumerate}[resume]
  \item Kehilangan \qty{1}{L} akan, tanpa pengaturan, menurunkan
        $\bar P$ sebesar \qty{45}{mmHg}. Hitunglah tekanannya setelah
        refleksnya.
  \item Refleksnya menyempitkan \omterm{def:b2:blood-circulation:vessels}{arteriol} kulit dan ususnya sehingga
        hambatannya melipat dua. Hitung ulang hambatan totalnya dan
        aliran ke kulit, usus, dan otaknya pada tekanan pada
        pertanyaan 13.
  \item Selama sejam berikutnya, \qty{500}{mL} cairan antarsel masuk
        ke dalam \omterm{def:b2:blood-circulation:vessels}{kapilernya}. Jelaskan mekanismenya dengan \omterm{prop:b2:blood-circulation:exchange}{gaya
        Starling}, dan apa yang terjadi pada \omterm{def:b2:blood-circulation:blood}{hematokritnya}.
  \item Selama tiga hari ginjalnya memulihkan volume \omterm{def:b2:blood-circulation:blood}{plasmanya}.
        Sebutkan kedua hormon yang bertanggung jawab beserta
        pemicunya.
  \item Berapa lama sumsumnya memerlukan waktu untuk mengganti \omterm{def:b2:blood-circulation:blood}{sel
        darah merahnya} pada \num{2.5} juta tiap detik (\num{5e12} tiap
        liter)?
  \item Kehilangan \qty{2.5}{L} menjenuhkan refleksnya. Jelaskan,
        dengan batinya dan kurva refleksnya, mengapa tekanannya lalu
        runtuh.
\end{enumerate}

\medskip
\noindent\textbf{Bagian IV --- Berlari.}
Pada olahraga yang berat hambatan ototnya turun menjadi
\qty{0.33}{mmHg.s/mL}, kulitnya menjadi 5, usus dan ginjalnya melipat
dua, hambatan otak dan \omterm{def:b2:heart:anatomy}{jantungnya} turun sehingga otaknya menjaga
alirannya dan otot \omterm{def:b2:heart:anatomy}{jantungnya} melipatempatkan alirannya pada tekanan
yang baru; $\bar P$ naik menjadi \qty{110}{mmHg}.
\begin{enumerate}[resume]
  \item Hitunglah aliran ke otot dan ke kulitnya.
  \item Hitunglah hambatan totalnya dan \omterm{thm:b2:heart:starling}{curah jantungnya}.
  \item Pada laju 180, hitunglah \omterm{prop:b2:heart:cycle}{volume sekuncupnya}.
  \item Oksigen \omterm{def:b2:blood-circulation:vessels}{venanya} turun menjadi \qty{40}{mL/L}. Hitunglah
        serapan oksigennya dan faktor yang dengannya ia melampaui
        keadaan istirahatnya.
  \item Bagikan faktor itu kepada alirannya dan kepada
        penyaripatiannya.
  \item Jelaskan bagaimana \omterm{def:b2:blood-pressure:baroreflex}{barorefleksnya} dapat mengizinkan tekanan
        \qty{110}{mmHg} alih-alih membetulkannya.
  \item Nyatakan hasilnya: penurunan sisa ketika berdiri, tekanan
        setelah pendarahan \qty{1}{L}, serta keluaran dan serapan
        oksigennya ketika berolahraga.
\end{enumerate}
\end{problem}
